% Mathématiques 5e — cours V3
% Pensée informatique et algorithmique par blocs

\section{Objectifs}

\begin{itemize}
\item Manipuler des instructions simples et les séquencer.
\item Identifier les entrées et les sorties d'un programme.
\item Représenter des formules par une expression informatique dans un langage de programmation par blocs.
\item Calculer la valeur d'une formule à l'aide d'une suite d'instructions.
\item Prévoir la valeur d'une expression informatique avant son exécution.
\item Analyser un programme simple donné et modifier ses paramètres.
\item Effectuer une boucle inconditionnelle répétant une séquence un nombre fixé de fois.
\end{itemize}

\section{1. Une instruction précise}

Une instruction décrit une action élémentaire.

Par exemple :
\begin{itemize}
\item lire une valeur ;
\item ajouter $3$ ;
\item multiplier par $2$ ;
\item avancer de $10$ pas ;
\item afficher un résultat.
\end{itemize}

Une instruction doit être suffisamment précise pour être exécutée sans interprétation ambiguë.

\section{2. Séquence d'instructions}

Une séquence est une suite ordonnée d'instructions.

\begin{exemple}
Programme :
\begin{verbatim}
lire un nombre
ajouter 4
multiplier par 3
afficher le résultat
\end{verbatim}

Avec l'entrée :
\[
2,
\]
on obtient :
\[
2+4=6,
\]
puis :
\[
6\times3=18.
\]

La sortie vaut :
\[
\boxed{18}.
\]
\end{exemple}

L'ordre des instructions change le résultat.

\section{3. Entrée et sortie}

\begin{definition}
Une \textbf{entrée} est une donnée fournie au programme.

Une \textbf{sortie} est une information produite par le programme.
\end{definition}

\coursefigure{/images/mathematiques/5eme/informatique-blocs.svg}{Suite de blocs représentant une entrée, un calcul et une sortie.}{Un programme simple peut être lu comme une chaîne : donnée d'entrée, transformation, résultat de sortie.}

\section{4. Exemple d'entrée numérique}

Une procédure reçoit :
\[
x=5.
\]

Elle calcule :
\[
2x+3.
\]

La sortie vaut :
\[
2\times5+3=13.
\]

Donc :
\[
\boxed{13}.
\]

La valeur de $x$ est l'entrée ; le nombre $13$ est la sortie.

\section{5. Variable lue en entrée}

En 5e, la notion de variable informatique est abordée surtout comme une donnée que le programme lit.

On peut imaginer un bloc :
\begin{verbatim}
demander un nombre
mettre la réponse dans x
\end{verbatim}

Le programme utilise ensuite la valeur mémorisée.

\begin{attention}
L'objectif n'est pas encore de manipuler en autonomie des variables complexes ou des structures conditionnelles : ces notions sont approfondies en 4e.
\end{attention}

\section{6. Traduire une formule}

Considérons :
\[
P=2L+2\ell.
\]

Un programme par blocs doit reproduire exactement cette formule.

Une séquence possible est :
\begin{verbatim}
lire L
lire l
calculer 2 × L
calculer 2 × l
additionner les deux résultats
afficher
\end{verbatim}

L'informatique doit respecter les mêmes priorités que le calcul mathématique.

\section{7. Expression informatique}

Une formule mathématique peut être écrite sous une forme proche d'un langage de programmation.

Par exemple :
\[
A=\pi r^2
\]
peut être représentée par :
\begin{verbatim}
pi * r * r
\end{verbatim}

ou, selon l'environnement :
\begin{verbatim}
pi * (r ^ 2)
\end{verbatim}

Il faut connaître les conventions du langage utilisé.

\section{8. Prévoir avant d'exécuter}

Avant de lancer un programme, il est utile de prévoir le résultat.

Considérons :
\begin{verbatim}
x = 4
résultat = x * x + 3
\end{verbatim}

On prévoit :
\[
4^2+3=16+3=19.
\]

Si le programme affiche autre chose, il faut chercher une erreur.

\section{9. Table de suivi}

Une table permet de suivre les valeurs intermédiaires.

Pour le programme :
\begin{verbatim}
lire x
ajouter 2
doubler
soustraire 1
\end{verbatim}

avec :
\[
x=3,
\]
on obtient :

\begin{tabular}{c|c}
Étape & Valeur \\
entrée & $3$ \\
après $+2$ & $5$ \\
après le double & $10$ \\
après $-1$ & $9$ \\
\end{tabular}

La sortie vaut :
\[
\boxed{9}.
\]

\section{10. Modifier un paramètre}

Un paramètre est une valeur que l'on peut modifier sans changer toute la structure du programme.

Par exemple :
\begin{verbatim}
avancer de 50
\end{verbatim}

Le nombre :
\[
50
\]
est un paramètre.

Le remplacer par :
\[
80
\]
allonge le déplacement tout en conservant la même instruction.

\section{11. Analyser un programme fourni}

Pour comprendre un programme, on peut se demander :
\begin{itemize}
\item quelles sont les entrées ?
\item quelles opérations sont exécutées ?
\item dans quel ordre ?
\item quels paramètres peuvent être modifiés ?
\item quelle est la sortie ?
\end{itemize}

Cette lecture structurée évite de modifier le programme au hasard.

\section{12. Exemple de programme géométrique}

Un programme de dessin contient :
\begin{verbatim}
avancer de 80
tourner de 90 degrés
avancer de 80
tourner de 90 degrés
\end{verbatim}

On peut prévoir le trajet avant l'exécution.

La valeur :
\[
80
\]
contrôle la longueur.

La valeur :
\[
90^\circ
\]
contrôle le changement de direction.

\section{13. Répéter une séquence}

Lorsque les mêmes instructions doivent être exécutées plusieurs fois, on utilise une boucle.

\coursefigure{/images/mathematiques/5eme/informatique-boucle.svg}{Bloc répéter cinq fois contenant deux instructions.}{Une boucle inconditionnelle répète une séquence un nombre fixé de fois sans tester de condition.}

\begin{definition}
Une \textbf{boucle inconditionnelle} répète une séquence un nombre de fois fixé à l'avance.
\end{definition}

\section{14. Exemple de boucle}

Pour tracer un carré :
\begin{verbatim}
répéter 4 fois
    avancer de 60
    tourner de 90 degrés
\end{verbatim}

À chaque répétition, les deux instructions sont exécutées.

Il y a :
\[
4
\]
déplacements et :
\[
4
\]
rotations.

\section{15. Nombre total d'instructions}

Si une boucle répète :
\[
5
\]
fois une séquence contenant :
\[
3
\]
instructions, le nombre total d'exécutions élémentaires vaut :
\[
5\times3=15.
\]

Ce calcul peut aider à prévoir la durée ou la structure du programme.

\section{16. Boucle et formule}

Une boucle peut aussi servir à répéter un calcul.

\begin{verbatim}
répéter 6 fois
    ajouter 3 au total
\end{verbatim}

Si le total vaut initialement :
\[
0,
\]
la valeur finale est :
\[
6\times3=18.
\]

La boucle traduit ici une multiplication répétée.

\section{17. Modifier le nombre de répétitions}

Dans :
\begin{verbatim}
répéter 4 fois
    avancer de 50
    tourner de 90 degrés
\end{verbatim}

si l'on remplace :
\[
4
\]
par :
\[
3,
\]
le dessin n'est plus un carré complet.

Modifier un seul paramètre peut donc changer fortement le comportement.

\section{18. Modifier un angle}

Pour un polygone régulier construit par déplacement, l'angle de rotation joue un rôle essentiel.

Une valeur incorrecte produit une figure qui ne se referme pas.

L'élève peut :
\begin{itemize}
\item prévoir l'effet ;
\item tester ;
\item comparer au résultat attendu ;
\item corriger le paramètre.
\end{itemize}

\section{19. Débogage}

Déboguer consiste à rechercher puis corriger une erreur.

\begin{methode}
\begin{enumerate}
\item prévoir le résultat attendu ;
\item exécuter ou simuler pas à pas ;
\item repérer la première étape incorrecte ;
\item identifier l'instruction ou le paramètre responsable ;
\item corriger ;
\item tester à nouveau.
\end{enumerate}
\end{methode}

\section{20. Exemple de débogage numérique}

On veut calculer :
\[
3x+2
\]
pour :
\[
x=4.
\]

Le programme produit :
\[
18.
\]

Or le calcul attendu est :
\[
3\times4+2=14.
\]

Si le programme contient :
\begin{verbatim}
multiplier x par 4
ajouter 2
\end{verbatim}
l'erreur se situe dans le paramètre :
\[
4
\]
qui devrait être :
\[
3.
\]

\section{21. Exemple de lecture de formule}

Le programme :
\begin{verbatim}
lire x
mettre résultat à x * x
ajouter 5 à résultat
afficher résultat
\end{verbatim}

traduit :
\[
x^2+5.
\]

Pour :
\[
x=3,
\]
on prévoit :
\[
3^2+5=14.
\]

\section{22. Expression et priorités}

Dans une expression informatique :
\begin{verbatim}
2 + 3 * 4
\end{verbatim}
la multiplication est généralement prioritaire.

Le résultat attendu est :
\[
2+12=14.
\]

Pour obtenir :
\[
20,
\]
il faudrait écrire :
\begin{verbatim}
(2 + 3) * 4
\end{verbatim}

Les parenthèses conservent donc leur rôle mathématique.

\section{23. Programme et tableur}

Un tableur et un programme par blocs peuvent calculer une même formule.

Par exemple :
\[
y=2x+1.
\]

Un tableur peut produire automatiquement plusieurs valeurs.

Un programme par blocs peut demander une valeur de $x$ puis afficher la valeur de $y$.

Les outils diffèrent, mais la relation mathématique reste identique.

\section{24. Méthode de lecture}

\begin{methode}
Pour analyser un programme :
\begin{enumerate}
\item identifier les entrées ;
\item lire les instructions dans l'ordre ;
\item repérer les calculs ;
\item repérer les boucles ;
\item noter les paramètres ;
\item prévoir la sortie ;
\item exécuter et comparer.
\end{enumerate}
\end{methode}

\section{25. Limite du programme de 5e}

\begin{attention}
En 5e, on utilise une variable principalement comme donnée saisie et on travaille la boucle à nombre de répétitions fixé.

Les instructions conditionnelles et une manipulation plus autonome des variables sont introduites en 4e.
\end{attention}

Cette limite évite de mélanger les niveaux du programme.

\section{26. Erreurs fréquentes}

\begin{itemize}
\item Confondre entrée et sortie.
\item Exécuter les instructions dans le mauvais ordre.
\item Oublier une priorité opératoire lors de la traduction d'une formule.
\item Modifier plusieurs paramètres à la fois lors d'un débogage.
\item Confondre nombre de répétitions et nombre d'instructions dans la boucle.
\item Oublier qu'une boucle exécute tout son contenu à chaque tour.
\item Introduire des conditions alors que la boucle demandée est inconditionnelle.
\item Croire que la machine garantit la justesse de la formule programmée.
\end{itemize}

\section{À retenir}

\begin{aretenir}
Un programme reçoit éventuellement des \textbf{entrées}, exécute une \textbf{séquence d'instructions} et produit une \textbf{sortie}.

Une formule mathématique peut être traduite en expression informatique.

Il faut savoir prévoir le résultat avant l'exécution.

Une boucle inconditionnelle répète une séquence un nombre fixé de fois :
\[
\boxed{\text{répéter }n\text{ fois}}.
\]

En 5e, l'essentiel est de savoir lire, prévoir, modifier et vérifier un programme simple plutôt que d'écrire des programmes complexes en autonomie.
\end{aretenir}
